\scanpage{36}
équisingularité) et fermé (par propreté) donc \uline{égal} à $\Phi$ qui est connexe; \deleted{et d'autre part} donc $\Phi-f(F)$ est une réunion d'images de feuilles \inserted{$(\subset\overline F\cap f^{-1}(\Phi))$} de $X_i$, donc ouvert dans $\Phi$, donc $f(F)$ est à la fois ouvert et fermé dans $\Phi$, donc $=\Phi$ ] De plus, \deleted{bien entendu} comme $F$ est fibré loc. trivial sur l'ouvert $f(F)$ de $\Phi$, \deleted{\unsure{et que}} (équisingularité) il est en fait fibré loc. trivial sur $\Phi$ tt entier. Mieux, $f^{-1}(\Phi)$ \deleted{est fibré} \uline{avec} \deleted{la}\inserted{sa} \uline{partition induite} est un fibré loc. trivial sur $\Phi$. Comme $F$ \deleted{et} et $\Phi$ sont des variétés, il s'ensuit que les fibres de $F$ sur $\Phi$ doivent aussi être des variétés ---

Dans la première lettre à A-Campo, j'ai donné une description conjecturale de la ``partition privilégiée canonique'' d'un espace $X$, en termes de points ``immobiles''. Voici une description conjecturale de la ``partition privilégiée'' d'une application \uline{propre} \deleted{$f$}
\[
f:X\to Y
\]
Pour simplifier, on suppose $Y$ compact, donc $X$ compact. On \deleted{\unsure{détermine par récurrence une}}\inserted{procède par récurrence sur $\dim Y$ pour} \deleted{filtration décroissante de} la construction. Considérons le plus grand ouvert $U$ de $Y$ tel que 1°) $U$ est non singulier, 2°) $X|U\unsure{-|\to}U$ est un fibré loc. trivial. \unsure{Il faut} prouver que $Y-U=Y'$ est un sous-espace (de la catégorie ``par morceaux'' envisagée) \uline{rare} (donc $\dim Y'<\dim Y$). Les composantes connexes de $Y-Y'$ seront les feuilles
\mspage{37}
